\section{Conclusion and Future Work}
\label{sec:conclusion}

This paper presented a multi-stage framework for PE malware detection that
connects evidence fusion, explanation-guided model adaptation, and quantitative
explanation evaluation. The framework uses a LightGBM static classifier and
native TreeSHAP, defines a stateful PID rule for aligned static--dynamic scores,
selects active-learning queries through predictive uncertainty, directional
entropy, and signed SHAP variance, and tests feature rankings with controlled
deletion and equal-budget random masks.

The verified evidence supports two conclusions within a row-disjoint internal
pilot. First, the full acquisition policy improved ranking at a fixed 700-label
training budget: mean ROC--AUC was 0.9749, compared with 0.9691 for uncertainty
sampling and 0.9634 for random selection. The full policy exceeded uncertainty
sampling in AUC for all five paired seeds. Its FPR was lower in only two seeds,
so the result does not establish stable false-positive reduction. Second,
top-$|\mathrm{SHAP}|$ deletion affected the fixed fidelity model much more than
random deletion. At $k=25$, AUC fell from 0.9802 to 0.8185, while random masks
retained mean AUC 0.9799. This supports the relevance of the SHAP ranking to model
behavior under training-mean replacement, not causal or executable-level
validity.

The paper makes no empirical PID performance claim. The repository lacks real
dynamic observations aligned to the static validation/test rows, and fusion
replay time would not measure the cost of acquiring such observations. The
controller, validation grid, state policy, baselines, and factorial evaluation
are specified so that a subsequent experiment can test rather than assume its
benefit.

Three priorities govern future work. First, a dataset manifest must establish
file identities, collection time, source, family, license, and label semantics;
this enables temporal, family-held-out, and external validation. Second, dynamic
reports should be collected under a locked observation policy and paired to the
static rows, followed by Static/Dynamic/Average/Fixed/P/PI/PID ablations and
separate latency decomposition. Third, fidelity evaluation should move from
coordinate-wise means toward schema-aware group interventions, manifold checks,
and, where possible, functionality-preserving PE transformations. Further work
should also compare diversity-aware acquisition, multiple fitted fidelity models,
analyst annotation cost, and robustness to poisoning or explanation manipulation.

The broader methodological lesson is that a multi-stage security system should
not be validated by one headline AUC. Its classifier, acquisition policy,
explanations, evidence fusion, and operational timing answer different questions.
Keeping those claims separate---and attaching each to replayable evidence---is a
necessary condition for credible progress from a thesis prototype toward a
deployable detection workflow.
